problem:  
Let \((M^n,g,\psi)\) be a compact Riemannian spin manifold admitting a **real Killing spinor** with constant \(\lambda \neq 0\):  
\[
\nabla_X^g\psi = \lambda\, X\!\cdot\!\psi, \qquad \forall X \in TM,
\]
so that \(g\) is Einstein with \(\mathrm{Ric}(g) = 4\lambda^2(n-1)g\).  
We consider infinitesimal variations of the pair \((g,\psi)\):
\[
g_t = g + t\,h + O(t^2), \qquad \psi_t = \psi + t\,\dot\psi + O(t^2),\qquad \lambda_t = \lambda + t\,\dot\lambda + O(t^2),
\]
where \(h\in \Gamma(\mathrm{Sym}^2 T^*M)\) and \(\dot\psi \in \Gamma(S)\).

---

**Step 1: Linearized equations.**  
The infinitesimal form of the Killing spinor equation is obtained by differentiating
\(\nabla^{g_t}_X \psi_t = \lambda_t\, X\!\cdot_{g_t}\! \psi_t\):
\[
\nabla_X^g \dot\psi
+ \tfrac{1}{4}(\nabla_X h)_{ab}\, e^a\!\cdot\! e^b\!\cdot\!\psi
- \lambda\, X\!\cdot\!\dot\psi
- \dot\lambda\, X\!\cdot\!\psi = 0, \qquad \forall X,
\]
where \(\{e_a\}\) is a local orthonormal frame, and \((\nabla_X h)_{ab}\) is contracted via Clifford multiplication.

Define the **linearized Killing-spinor operator**
\[
L_{(g,\psi)}(h,\dot\psi) := \nabla\dot\psi - \lambda\,(\cdot)\!\cdot\!\dot\psi - \tfrac{1}{4}(\nabla(\cdot) h)_{ab}\, e^a\!\cdot\! e^b\!\cdot\!\psi.
\]
Then \((h,\dot\psi)\) solves the linearized equation iff \(L_{(g,\psi)}(h,\dot\psi) + \dot\lambda\,(\cdot)\!\cdot\!\psi = 0\).

---

**Step 2: Diffeomorphism invariance and gauge fixing.**  
The action of an infinitesimal diffeomorphism \(X \in \Gamma(TM)\) on the pair \((g,\psi)\) is
\[
\delta_X g = L_X g, \qquad
\delta_X \psi = \nabla_X \psi + \tfrac{1}{4}(dX^\flat)\!\cdot\!\psi.
\]
To fix this gauge, impose the **Bianchi condition**
\[
\delta h + \tfrac{1}{2}d(\mathrm{tr}\,h)=0,
\]
ensuring ellipticity of the linearized Einstein operator
\[
E_g(h) = \tfrac{1}{2}\Delta_L h - 2\lambda^2(n-1)\,h,
\]
where \(\Delta_L\) is the Lichnerowicz Laplacian.

---

**Step 3: Coupled operator and complex.**  
Define the combined operator
\[
\mathfrak{P}_{(g,\psi)} =
(E_g, L_{(g,\psi)}):
\Gamma(\mathrm{Sym}^2 T^*M)\oplus \Gamma(S)
\longrightarrow
\Gamma(\mathrm{Sym}^2 T^*M)\oplus \Gamma(T^*M\!\otimes\! S).
\]
Consider the infinitesimal complex
\[
0 \longrightarrow \Gamma(TM)
\xrightarrow{\;\mathcal{L}_{(g,\psi)}\;}
\Gamma(\mathrm{Sym}^2 T^*M)\oplus \Gamma(S)
\xrightarrow{\;\mathfrak{P}_{(g,\psi)}\;}
\Gamma(\mathrm{Sym}^2 T^*M)\oplus \Gamma(T^*M\!\otimes\! S)
\longrightarrow 0,
\]
where \(\mathcal{L}_{(g,\psi)}(X) = (L_X g,\, \nabla_X \psi + \tfrac{1}{4}(dX^\flat)\!\cdot\!\psi)\).

---

**Question:**  
Determine whether the complex above is **elliptic** in the sense of Atiyah–Singer, i.e. whether the induced symbol sequence
\[
0 \longrightarrow (T_xM) \xrightarrow{\;\sigma_\xi(\mathcal{L})\;}
\mathrm{Sym}^2(T_x^*M)\oplus S_x
\xrightarrow{\;\sigma_\xi(\mathfrak{P})\;}
\mathrm{Sym}^2(T_x^*M)\oplus (T_x^*M\!\otimes\!S_x)
\longrightarrow 0,
\]
is exact for all nonzero covectors \(\xi\).  

If ellipticity fails generally, characterize its geometric obstruction—does failure correspond to specific **Lichnerowicz, Dirac, or cone spectral resonances**?  
Conversely, if the complex is elliptic for certain classes (e.g. 6D nearly Kähler, 7D nearly parallel \(G_2\)), can one describe the corresponding **finite-dimensional moduli** space of Killing-spinor-preserving Einstein deformations?

---

Why is it a "good" problem:  
This problem isolates a fundamental analytic question: *does Killing-spinor geometry admit an elliptic deformation theory parallel to the special-holonomy case?* It is compelling because:

- It builds a concrete **coupled geometric-analytic complex** intertwining the Einstein operator and the spinorial Killing equation—making it accessible yet profound.  
- Establishing ellipticity would yield an analytic foundation for moduli of Einstein metrics with real Killing spinors, analogous to the deformation complexes for Calabi–Yau or \(G_2\) structures.  
- Failure of ellipticity would identify deep spectral or holonomy obstructions linked to the cone’s geometry, revealing structural constraints connecting spinorial fields and curvature.  
- The formulation bridges **spin geometry**, **elliptic operator theory**, **Einstein deformation analysis**, and **holonomy theory**, offering a genuinely interdisciplinary research direction.  

Its elegant simplicity—“is there an elliptic deformation complex for Killing spinor geometry?”—conceals significant analytical and geometric depth, fulfilling all the hallmarks of a *“good” research-level mathematical problem*.